
\section{The Turns Idea.}

The main problem we have is that, for now, the history is defined only for configurations that, at time $t$, see a non-trivial syndrome. What we can do is imagine that at each turn we flip the bits according to a fixed order. We think of the basic time unit as a turn, and each real-time iteration contains $t n^{d}$ turns. Now, if a bulk of bits $A$ is flipped at time $t$, then there must be a turn $\tau \le t n^{d}$ such that at that turn there was a configuration that has a history.

Now we would like to describe what the history of that configuration looks like, and we do that by connecting the turn $\tau$ to the closest $t^{\prime}n^{d}$. The idea is that we cann't ensure the shrinking lemma. Yet we can tell for sure that if there is expansion then its limited.  


I have also the felling that the shrinking works, and the idea goes as follow, assume that $v \in \Gamma^{(t-1)}$ gets the maximum of $L_{i}$ either we stop before flipping a bit that is 'supported' on $v$ or not, in the first case, we get that any $u$ that is added by $v^{\prime}$ might has $L_{i}(u) = L_{i}(v^{\prime}) < L_{i}(v)$, in the second case we have already handled. 

The turn machinery makes sense only in the finite case. Yet, we can back to the Tree Separator Section, and show that with probability  
\begin{equation*}
  \begin{split}
    n^{d} \cdot  t n^{d} q^{n}
  \end{split}
\end{equation*}
We don't have a big trees which implies independence. Then we can show that with probability:  
\begin{equation*}
  \begin{split}
    n^{d}t (n^{d})^{2}q^{\prime \delta n}
  \end{split}
\end{equation*}
There was no turn in which it had a configuration that was connected at the time, saw a syndrome, and had a side-length that behaved like $\Theta(n)$.

\ctt{ The model has to redefined again such it will be clear that outcomes are possibles running of the system. Currently this can not be inferred. }



\newcommand{\comDAG}{\pi[C,\mathcal{E}]} 
Informally, a quantum circuit is a placement of a collection of quantum gates, taken from a finite set, in different space-time locations. Each space-time location, associated with a gate, encodes the time of the application and the qubits in the support of the gate. Formally
\begin{definition}[Quantum Circuit]
A space-time location $l$ is a tuple that encodes in its first coordinate time, and in its second coordinate a subset of integers referring to a support of (q)bits. For a collection of quantum gates $\mathcal{G} = \{ g \colon L(\mathcal{H}_{2}^{\otimes \Delta} ) \rightarrow L(\mathcal{H}_{2}^{\otimes \Delta} ) \}$, a quantum circuit $C$, with depth $d$ and width $w$, is a collection of tuples $\{ u_{i} = (g_{i},l_{i}) \}$, where $u_{i}$ are the gates of $C$ and each $u_{i}$ is composed of a gate $g_{i} \in \mathcal{G}$ and a space-time location $l_{i} \in \mathbb{Z}_{d} \times \mathbb{Z}_{w}^{\Delta}$. The gates of $C$ have to satisfy the constraint that for every two gates with the same time coordinate, there cannot be a non-empty intersection in their space coordinate, namely for $u_{i}=(g_{i},l_{i}),u_{j}=(g_{j},l_{j}) \in C$, the equality $l_{i,1} = l_{j,1}$ implies $l_{i,2} \cap l_{j,2} = \emptyset$.


\end{definition}



\begin{definition}[Associated Computation DAG]
  Let $C$ be a quantum circuit, with depth $d$ and width $w$. For a set of faults location $\mathcal{E} \subset [d] \times [w] $ we define the associated computation DAG of $(C,\mathcal{E})$ to be the acyiclic directed graph $\comDAG = (\mathcal{U}, E)$ as follow: 
  \begin{enumerate}
    \item The vertices set $\mathcal{U}$ equals $G[C] \cup \{ N_{\text{loc} } \colon \text{loc } \in \mathcal{E}$, namely the union of the gates of $C$ and the faults in $\mathcal{E}$. 
      \item 
  \end{enumerate}
\end{definition}


\begin{definition}[Unfolded configuration]  
  For a quantum circuit $C$, and a set of faults $\mathcal{E}$, Let $\comDAG = (\mathcal{U},E)$ be their associated computation DAG. Let $f$
  For a running configuration $\omega \in \Omega$ we define the unfolded-configuration. 
\end{definition}




\begin{claim}
  Let $\omega \in \Omega$ be an outcome, 
\end{claim}
